\section{The fixed-shift arithmetic boundary}
\label{sec:fixed-shift-boundary}

The completed signal can recover every finite coefficient, but recovery
is not evaluation.  Two exact slices locate the remaining arithmetic
difficulty.

\subsection{The endpoint prime-pair slice}

The factor endpoint $m=1$ is characterized in primitive coordinates by
\[
 (a,b,k)=(n,1,1).
\]
Define its fully diagonal energy, independently of the source cutoff, by
\begin{equation}
 \cF_h^{\rm end}(X)
 :=\sum_{n\geq1}
 (\Lambda(n)-1)^2(\Lambda(n+h)-1)^2
 |W(n/X)|^2.
 \label{eq:endpoint-energy}
\end{equation}
For $D\leq X^{1-\eta}$ this is exactly the $m=1$ slice of
$\cF_{h,D}(X)$ once $X$ is large, because $n\geq\alpha X>D$ on the
support.

\begin{proposition}[Endpoint reduction]
\label{prop:endpoint-reduction}
For every fixed nonzero integer $h$,
\begin{equation}
 \cF_h^{\rm end}(X)
 =(\log X)^2
 \sum_n\Lambda(n)\Lambda(n+h)|W(n/X)|^2
 +O_{h,W}(X\log X).
 \label{eq:endpoint-reduction}
\end{equation}
\end{proposition}

\begin{proof}
On the support, $n,n+h\asymp_WX$.  If both are prime, then
\[
 (\log n-1)^2(\log(n+h)-1)^2
 =(\log X)^2\log n\log(n+h)+O_W((\log X)^3).
\]
The Selberg upper-bound sieve gives
$O_{h,W}(X/(\log X)^2)$ simultaneous prime values, so their total error
is $O(X\log X)$.  If exactly one of the two integers is prime and the
other is not a prime power, the summand on the left is
$O((\log X)^2)$; the prime number theorem bounds the total by
$O(X\log X)$.  If neither is a prime power, the total is $O(X)$.

Finally, the number of proper prime powers in an interval of size
$O(X)$ is $O(\sqrt X)$ up to a harmless logarithmic factor.  Giving all
terms involving one such value the trivial weight $O((\log X)^4)$
produces $o(X\log X)$.  These cases also cover the proper-prime-power
part of the von Mangoldt correlation on the right.
\end{proof}

For a fixed nonzero shift, let
\begin{equation}
 \mathfrak S(h):=
 \begin{cases}
 2C_2\displaystyle\prod_{\substack{p\mid h\\p>2}}
 \dfrac{p-1}{p-2},&2\mid h,\\[0.8em]
 0,&2\nmid h,
 \end{cases}
 \qquad
 C_2:=\prod_{p>2}\left(1-\frac1{(p-1)^2}\right).
 \label{eq:prime-pair-singular-series}
\end{equation}
This is the Hardy--Littlewood prime-pair singular series
\citep{HardyLittlewood1923}; it is not assumed in the positive theorems.

\begin{corollary}[Exact Hardy--Littlewood strength threshold]
\label{cor:endpoint-HL-equivalence}
For every fixed even $h\ne0$, the asymptotic
\begin{equation}
 \cF_h^{\rm end}(X)
 =\mathfrak S(h)I_2(W)X(\log X)^2
 +o_{h,W}(X(\log X)^2)
 \label{eq:endpoint-predicted-asymptotic}
\end{equation}
holds if and only if
\begin{equation}
 \sum_n\Lambda(n)\Lambda(n+h)|W(n/X)|^2
 =\mathfrak S(h)I_2(W)X+o_{h,W}(X).
 \label{eq:weighted-HL-correlation}
\end{equation}
For $h=2$, this is the corresponding weighted twin-prime
Hardy--Littlewood asymptotic.
\end{corollary}

\begin{proof}
Subtract \eqref{eq:endpoint-reduction} and divide by $(\log X)^2$.
Its error becomes $O(X/\log X)=o(X)$, proving both implications.
\end{proof}

When $W\not\equiv0$, \eqref{eq:weighted-HL-correlation} at $h=2$ would
in particular imply infinitely many twin primes after proper prime
powers are removed.  Neither side of the equivalence is proved here.
For odd fixed $h$, elementary parity and the same case split as above
give only the much smaller statement
\begin{equation}
 \cF_h^{\rm end}(X)=O_{h,W}(X\log X).
\label{eq:odd-endpoint-bound}
\end{equation}
Indeed, a fixed odd difference excludes two prime values at height
$X$; proper prime powers are negligible, and the terms with exactly one
prime have total size $O(X\log X)$.

\subsection{A fixed quadratic ray persists}

There is an earlier obstruction before the endpoint correlation.  On
the primitive ray $(a,b)=(6,1)$ and at $h=1$, the source value satisfies
$u(6k)=-1$ for every $k\geq1$.  Define the cutoff-free comparison
\begin{equation}
 Q_6(X):=
 \sum_{k\geq1}(\Lambda(6k^2+1)-1)^2|W(6k^2/X)|^2.
\label{eq:quadratic-ray-energy}
\end{equation}
If $D<\sqrt{6\alpha X}$, then every supported $k$ satisfies $6k>D$, so
$Q_6(X)$ is exactly the $(6,1)$-ray slice of $\cF_{1,D}(X)$.  For larger
$D$ we make no such identification.
Put
\begin{equation}
 P_6(X):=
 \sum_{k\geq1}\Lambda(6k^2+1)|W(6k^2/X)|^2.
 \label{eq:quadratic-prime-sum}
\end{equation}

\begin{proposition}[Quadratic-prime strength threshold]
\label{prop:quadratic-prime-threshold}
One has
\begin{equation}
 Q_6(X)=\log X\,P_6(X)+O_W(\sqrt X).
 \label{eq:quadratic-ray-reduction}
\end{equation}
Consequently, for every fixed constant $C$,
\begin{equation}
 Q_6(X)=C\sqrt X\log X+o(\sqrt X\log X)
 \label{eq:quadratic-energy-asymptotic}
\end{equation}
if and only if
\begin{equation}
 P_6(X)=C\sqrt X+o(\sqrt X).
 \label{eq:quadratic-prime-asymptotic}
\end{equation}
\end{proposition}

\begin{proof}
Put $K\asymp_W\sqrt X$ for the supported $k$-range.  The
dimension-one Selberg upper-bound sieve for the irreducible quadratic
$6k^2+1$ gives
\[
 \#\{k\asymp_WK:6k^2+1\ {\rm is\ prime}\}
 \ll_W\frac K{\log K};
\]
see \citet{HalberstamRichert1974}.  Restoring the logarithmic
weight yields $P_6(X)\ll_WK\ll_W\sqrt X$.

For $r=1,2$, the proper-prime-power contribution satisfies
\[
 \sum_{\substack{k\asymp_W\sqrt X\\6k^2+1=p^j,\ j\geq2}}
 \Lambda(6k^2+1)^r|W(6k^2/X)|^2
 \ll_WX^{1/3}(\log X)^{r+1}+(\log X)^{r+1}
 =o_W(\sqrt X).
\]
For even $j$ this follows from the associated Pell equation, which has
$O_W(\log X)$ relevant solutions; for odd $j\geq3$, there are
$O_W(X^{1/3}\log X)$ candidates.  These estimates are detailed in
TPC-12 \citep{WangTPC12}.

At prime values in the support,
$\log(6k^2+1)=\log X+O_W(1)$.  Therefore
\[
 \begin{aligned}
 Q_6(X)
 &=
 \sum_k\Lambda(6k^2+1)^2|W(6k^2/X)|^2
 -2P_6(X)+\sum_k|W(6k^2/X)|^2\\
 &=\log X\,P_6(X)+O_W(P_6(X))+O_W(\sqrt X)
 =\log X\,P_6(X)+O_W(\sqrt X).
 \end{aligned}
\]
The final equivalence follows after division by $\log X$.
\end{proof}

The asymptotic \eqref{eq:quadratic-prime-asymptotic} is a weighted prime
number theorem for the fixed polynomial $6k^2+1$ and is open.  The radial
coordinate separates the individual $k$-labels but supplies no new
distribution theorem for their prime targets.

\subsection{What is lost after retaining only total energy}

The full two-coordinate signal is injective, but the scalar
$\cF_{h,D}=\sum|x_{a,b,k}|^2$ again forgets where its energy lies.  The
following coefficient-blind minimax identity, for arbitrary coefficient
vectors, prevents an endpoint conclusion from the total energy alone.

\begin{proposition}[Exact energy-only minimax boundary]
\label{prop:energy-minimax}
Let a finite index set $I$ be partitioned into nonempty sets $J$ and
$I\setminus J$.  For $x\in\C^I$, define
\[
 S(x)=\sum_{i\in I}|x_i|^2,
 \qquad
 Q_J(x)=\sum_{i\in J}|x_i|^2.
\]
Then, for every $B>0$,
\begin{equation}
 \inf_{\mathcal A:[0,B^2]\to\R}
 \sup_{\|x\|_2\leq B}
 |Q_J(x)-\mathcal A(S(x))|
 =\frac{B^2}{2}.
 \label{eq:energy-minimax}
\end{equation}
\end{proposition}

\begin{proof}
Take two vectors of norm $B$, one supported in $J$ and one in its
complement.  They have the same observed value $S=B^2$, while their
targets are $B^2$ and $0$.  Every estimator therefore has worst-case
error at least $B^2/2$.  Conversely, the estimator
$\mathcal A(S)=S/2$ obeys
$|Q_J-S/2|\leq S/2\leq B^2/2$.
\end{proof}

For the endpoint, $J$ is the set $b=k=1$.  The complete joint signal can
recover these coefficients by \eqref{eq:exact-coefficient-recovery}, but
the total-energy asymptotic \eqref{eq:short-shift-radial-energy} cannot
separate them from the bulk.  In particular,
$\cF_{h,D}\geq\cF_h^{\rm end}$ supplies no twin-prime lower bound: the
bulk may account for the entire known main term.
For $W\not\equiv0$ and large $X$, both the endpoint and its complement
are nonempty in the stated source-cutoff range.  The minimax statement
does not rule out an estimator using additional special arithmetic
information about the von Mangoldt coefficients.

The classical sieve parity limitation explains why upper-bound sieve
technology alone does not turn these reductions into lower bounds
\citep{Selberg1952Parity,Polymath2014}.  The present paper repairs a
geometric label collision and then stops exactly where new arithmetic
input is required.
